% -------------------------------------------------------------------------
% WBS ID: RES-NUM-ERR
% Deliverable: Numerical Error, Adaptivity and Quality Indicators
% Status: In progress
% Evidence status: Local conservation diagnostics and quality indicators established
% Review status: Not reviewed
% Last updated: 2026-07-18
% Dependencies: RES-NUM-MSH, RES-NUM-STB, RES-NUM-SPEC
% Downstream interfaces: RES-VER-CNV, RES-VER-MET, Software diagnostics
% -------------------------------------------------------------------------

\section{RES-NUM-ERR --- Numerical Error, Adaptivity and Quality Indicators}
\label{sec:res-num-err}

\subsection{Purpose and Scope}
\label{sec:res-num-err-purpose}

% TODO: Define what this Deliverable covers and explicitly excludes.

This Deliverable defines numerical diagnostics for the local
three-dimensional URANS--VOF model. It records what must be measured to
judge conservation, mesh quality, refinement needs and output
reliability. It does not set final validation tolerances.

\subsection{Research Questions}
\label{sec:res-num-err-research-questions}

% TODO: State the questions that must be answered before the Deliverable
% can be accepted.

\subsection{Terminology and Definitions}
\label{sec:res-num-err-definitions}

% TODO: Define the technical terminology, symbols, quantities and
% classifications used in this Deliverable.

\subsection{Literature Evidence}
\label{sec:res-num-err-literature-evidence}

% TODO: Record evidence from peer-reviewed papers, authoritative datasets,
% standards, technical reports and primary documentation.
%
% Do not create a detached literature-summary list. Explain how each source
% supports, contradicts or limits a project decision.

\textcite{roenby2016computational} and \textcite{vukcevic2018sharp}
are retained as evidence that VOF quality must be assessed through
phase conservation, boundedness and interface sharpness. Source role:
`3D-VOF`. Project conclusion: local diagnostic output shall include
both conservation quantities and engineering-output quality indicators.
Limitation: these diagnostics support numerical quality; they do not
replace physical validation against benchmark or experimental data.

\subsection{Governing Relations and Quantitative Definitions}
\label{sec:res-num-err-governing-relations}

% TODO: Record relevant equations, dimensionless groups, empirical models,
% extraction rules or quantitative definitions.
%
% Use align environments for displayed equations.
% Every displayed equation must have a unique \label{eq:...}.
% Do not place punctuation after displayed equations.

The global VOF-volume error over a timestep is:

\begin{align}
  E_{\alpha}^{n+1}
  &=
  \frac{
    M_{\alpha}^{n+1}
    -
    M_{\alpha}^{n}
    +
    \Delta t
    \sum_{f\in\partial\Omega}
    \phi_{\alpha,f}
  }{
    M_{\alpha}^{n}
    +
    \epsilon_{\alpha}
  }
  \label{eq:res-num-err-local-alpha-volume-error}
\end{align}

The global mass imbalance is:

\begin{align}
  E_m
  &=
  \frac{
    \left|
      \sum_i
      R_{m,i}
    \right|
  }{
    \sum_i
    \sum_{f\in\partial K_i}
    \left|
      \phi_f
    \right|
    +
    \epsilon_{\phi}
  }
  \label{eq:res-num-err-local-global-mass-imbalance}
\end{align}

\subsection{Discrete Formulation}
\label{sec:res-num-err-discrete-formulation}

% TODO: Complete this RES-NUM Domain-specific subsection.

Each accepted local timestep shall write a numerical-quality record
containing timestep size, Courant numbers, VOF boundedness corrections,
pressure residuals, mass residuals, turbulence residuals, mesh-quality
extrema and output-integral consistency checks.

\subsection{Conservation Properties}
\label{sec:res-num-err-conservation-properties}

% TODO: Complete this RES-NUM Domain-specific subsection.

The conservation diagnostics shall include:

\begin{itemize}
  \item global and local phase-volume balance;
  \item cellwise and global mass residuals;
  \item net boundary fluxes through inlet, outlet, side and top
    patches;
  \item momentum-flux balance across regional-to-local interfaces;
  \item force and moment balance on terrain and barrier patches where
    those quantities are extracted.
\end{itemize}

\subsection{Consistency and Formal Accuracy}
\label{sec:res-num-err-consistency-formal-accuracy}

% TODO: Complete this RES-NUM Domain-specific subsection.

\subsection{Stability Requirements}
\label{sec:res-num-err-stability-requirements}

% TODO: Complete this RES-NUM Domain-specific subsection.

\subsection{Mesh and Time-Step Requirements}
\label{sec:res-num-err-mesh-time-step-requirements}

% TODO: Complete this RES-NUM Domain-specific subsection.

The local refinement indicators shall include:

\begin{itemize}
  \item $\lVert\nabla\alpha\rVert$ near the air--water interface;
  \item local Courant and interface Courant numbers;
  \item pressure gradient and pressure impulse near the barrier;
  \item vorticity and turbulent-viscosity ratio in wakes and
    separation regions;
  \item wall-distance and wall-function validity indicators;
  \item force and moment sensitivity to mesh refinement;
  \item overtopping-volume sensitivity to crest and run-up resolution.
\end{itemize}

\subsection{Numerical Failure Modes}
\label{sec:res-num-err-numerical-failure-modes}

% TODO: Complete this RES-NUM Domain-specific subsection.

A simulation quality record shall flag:

\begin{itemize}
  \item clipped phase fractions or turbulence variables;
  \item repeated timestep rejection;
  \item pressure residual stagnation;
  \item large global mass imbalance;
  \item local mesh-quality excursions in impact or wall regions;
  \item abrupt pressure, force or moment spikes not supported by
    neighbouring mesh or timestep refinements;
  \item boundary-reflection indicators above tolerance.
\end{itemize}

\subsection{Verification Requirements}
\label{sec:res-num-err-verification-requirements}

% TODO: Complete this RES-NUM Domain-specific subsection.

Verification shall demonstrate that free-surface elevation, overtopping
volume, peak pressure, resultant force, impulse and moment converge
under mesh and timestep refinement before the local model is used to
compare candidate barriers.

\subsection{Candidate Approaches}
\label{sec:res-num-err-candidate-approaches}

% TODO: Define the credible candidate approaches considered.

\subsection{Critical Comparison}
\label{sec:res-num-err-critical-comparison}

% TODO: Compare candidates using physically and project-relevant criteria.
% Do not select an approach solely on popularity or implementation ease.

\subsection{Assumptions and Applicability}
\label{sec:res-num-err-assumptions}

% TODO: State assumptions and the conditions under which they are valid.

\subsection{Limitations and Failure Conditions}
\label{sec:res-num-err-limitations}

% TODO: State where the evidence, model or selected approach ceases to be
% reliable or applicable.

\subsection{Project-Specific Conclusions}
\label{sec:res-num-err-conclusions}

% TODO: Convert the evidence into conclusions specific to the Studentship
% project. Do not merely restate literature findings.

\subsection{Selected Approach and Rationale}
\label{sec:res-num-err-selected-approach}

% TODO: Record the selected approach only after sufficient evidence exists.
% State the alternatives rejected and the reasons for rejection.

The selected diagnostic approach is to report numerical quality
alongside every local engineering output. A force, moment or
overtopping value is not accepted as a standalone scalar unless the
supporting conservation, boundedness, mesh and timestep diagnostics are
available.

\subsection{Unresolved Decisions}
\label{sec:res-num-err-unresolved-decisions}

% TODO: Record unresolved questions, required evidence and decision owners.

Open diagnostic decisions are final numerical-quality thresholds,
adaptive-refinement triggers, output smoothing rules for pressure and
force histories, and the minimum convergence rate required between
mesh levels.

\subsection{Requirements and Handoffs}
\label{sec:res-num-err-requirements}

% TODO: State requirements passed to other Research Domains, Software,
% verification activities or later execution work.

Software shall emit the diagnostics listed in this Deliverable for
every local simulation used in validation, surrogate training or
candidate comparison. `RES-VER-MET` shall define which diagnostic
failures invalidate a physical validation comparison.

\subsection{Deliverable Completion Record}
\label{sec:res-num-err-completion-record}

% TODO: Record whether the Deliverable is:
%
% - Not started
% - In progress
% - Draft complete
% - Reviewed
% - Accepted
%
% Acceptance requires:
%
% - the research questions to be answered;
% - claims to be supported by appropriate evidence;
% - assumptions and limitations to be explicit;
% - the project-specific conclusion to be stated;
% - downstream requirements to be traceable;
% - unresolved decisions to be closed or formally deferred.
